\documentclass[10pt,a4paper]{amsart}
\usepackage[margin=19mm]{geometry}
\usepackage{amsmath,amssymb,mathtools}
\usepackage[scr=rsfs,cal=euler]{mathalfa}
\usepackage{xfrac}
\usepackage[unicode,hidelinks,bookmarks=false]{hyperref}
\newcommand{\ie}{i.e.\ }
\newcommand{\cf}{cf.\ }
\newcommand{\resp}{respectively\ }
\newcommand{\Subsection}{\subsection}
\providecommand{\nobreakdash}{\nobreak\hbox{-}}
\setlength{\emergencystretch}{2em}
\multlinegap0pt
\usepackage{amssymb}
\usepackage{float}
\newtheorem{lemma}{Lemma}
\newcommand{\lebesguespace}[1]{L^2(#1)}
\newcommand{\localvespace}{V^E_{h,k}}
\newcommand{\scalenergyproj}{\Pi^{\nabla}_{k,E}}
\title{Mass-Lumped Virtual Element Method with Strong Stability-Preserving Runge-Kutta Time Stepping for Two-Dimensional Parabolic Problems}
\author{Paulo Akira F. Enabe, Rodrigo Provasi}
\date{Source: arXiv:2510.06653v2 (CC BY 4.0)}
\begin{document}
\maketitle

\noindent\emph{Source and licence.} Mass-Lumped Virtual Element Method with Strong Stability-Preserving Runge-Kutta Time Stepping for Two-Dimensional Parabolic Problems. Version arXiv:2510.06653v2 (CC BY 4.0), distributed by arXiv under the Creative Commons Attribution 4.0 International licence, http://creativecommons.org/licenses/by/4.0/, as recorded in the arXiv licence metadata for this exact version and shown on the abstract page https://arxiv.org/abs/2510.06653v2. Full licence notice: \texttt{licenses/arxiv-2510.06653-CC-BY-4.0.txt}.\par\smallskip

\noindent\textit{Editorial scope.} Counted unit: the article's lemma eq:aux\_ineq\_energy\_stab (source0.tex lines 472--482). The article's complete original proof of it is delivered with it (source0.tex lines 484--527, one \texttt{proof} environment of 44 lines). No mathematics is written, repaired or re-derived: the statement and the proof are the article's own lines, in the article's own notation, and no retained line is substituted or re-flowed.  Retained uncounted as the source-local context the proof needs: lines 345--348; lines 463--465; lines 467--469; lines 593--595.  Nothing was omitted from any retained span, so the package carries no omission record.  No retained line contains a citation, so the wrapper carries no bibliography and no bibliography entry.  No retained line contains a figure environment, an image-inclusion command or drawing code, so no image is delivered and the package declares no asset dependency.  Editorial formatting only, in addition: an AMS \texttt{amsart} preamble that restates the article's own declarations and macros used by the retained lines, namely the environments \texttt{lemma} on separate counters, together with the standard packages those lines need.  The retained environments print in this document as Lemma 1 in order of appearance, whereas the article prints the same items with the numbering of its own full document.  The complete native source of the article as distributed in the arXiv e-print for this exact version is bundled unchanged as \texttt{sources/186624/source0.tex}.
\begin{equation}\label{eq:continuous_bilinear_stiff}
    \nonumber
    a^E(u,v) = \int \limits_E \nabla u \nabla v dE.
\end{equation}

\begin{equation}\label{eq:global_energy_bilinear_form}
    a_h(\cdot, \cdot) = \sum \limits_{E\in\mathcal{T}_h}a^E_h (\cdot, \cdot)
\end{equation}\

\begin{equation}\label{eq:global_mass_bilinear_form}
    m_h (\cdot, \cdot) = \sum \limits_{E\in \mathcal{T}_h} m^E_h (\cdot, \cdot) = \langle \cdot, \cdot \rangle_h.
\end{equation}

\begin{equation}\label{eq:stability_term_condition}
    C_0 a^E(w, w) \leq S^E(w,w) \leq C_1 a^E(w,w), \; \forall w \in \ker(\scalenergyproj),
\end{equation}

\begin{lemma}
    Consider the bilinear form in (\ref{eq:global_energy_bilinear_form}) and (\ref{eq:global_mass_bilinear_form}). There exist mesh independent constants $\alpha_*, \alpha^*, \beta_*, \beta^* > 0$ such that:
    \begin{equation}\label{eq:aux_ineq_energy_stab}
        \alpha_* \| \nabla v_h \|_{L^2(\Omega)}^2 \leq a_h (v_h, v_h) \leq \alpha^* \| \nabla v_h \|^2_{L^2(\Omega)}
    \end{equation}
    and
    \begin{equation}\label{eq:aux_ineq_mass_stab}
        \beta_* \| v_h \|_{L^2(\Omega)}^2 \leq \langle v_h, v_h \rangle_h \leq \beta^* \| v_h \|_{L^2{\Omega}}^2
    \end{equation}
    for all $v_h \in V_h$.
\end{lemma}

\begin{proof}
    Consider the decomposition for $v_h \in \localvespace$:
    \begin{equation}
    \nonumber
        v_h = \scalenergyproj v_h + w_h, \; \text{with} \; w_h = (I - \scalenergyproj)v_h \in \ker (\scalenergyproj).
    \end{equation}
    Considering the e local bilinear form in (\ref{eq:continuous_bilinear_stiff}), by orthogonality, it holds that:
    \begin{equation}
    \nonumber
        a^E(\scalenergyproj v_h, w_h) = 0.
    \end{equation}
    So, it is possible to write:
    \begin{equation}
    \nonumber
       \begin{split}
            \| \nabla v_h \|^2_{L^2(\Omega)} = a^E (\scalenergyproj v_h + w_h, \scalenergyproj v_h + w_h) = \\
            a^E(\scalenergyproj v_h, \scalenergyproj v_h) + 2 a^E(\scalenergyproj v_h, w_h) + a^E(w_h, w_h) = \\
            a^E(\scalenergyproj v_h, \scalenergyproj v_h) + a^E(w_h, w_h).
       \end{split}
    \end{equation}
    Thus,
    \begin{equation}
    \nonumber
        \| \nabla v_h \|^2_{L^2(\Omega)} = \| \nabla \scalenergyproj v_h \|^2_{L^2(\Omega)} + \| w_h \|^2_{L^2(\Omega)}.
    \end{equation}
    It holds true  that (this result is explored further in eq. (\ref{eq:stability_term_condition})):
    \begin{equation}
    \nonumber
        a^E_h(v_h, v_h) \geq \| \nabla \scalenergyproj v_h \|^2_{\lebesguespace{E}} + C_0 \| \nabla w_h \|^2_{\lebesguespace{E}} \geq \min (1, C_0) \| \nabla v_h \|^2_{\lebesguespace{E}},
    \end{equation}
    and
    \begin{equation}
    \nonumber
        a^E_h(v_h, v_h) \leq \| \nabla \scalenergyproj v_h \|^2_{\lebesguespace{E}} + C_1 \| \nabla w_h \|^2_{\lebesguespace{E}} \leq \max (1, C_1) \| \nabla v_h \|^2_{\lebesguespace{E}}.
    \end{equation}
    Summing over $E$ and using mesh regularity yields the global bounds with:
    \begin{equation}
    \nonumber
        \alpha_* = \min (1,C_0), \; \alpha^* = \max (1, C_1),
    \end{equation}
    resulting in (\ref{eq:aux_ineq_energy_stab}).

    Similar arguments and bounds hold for the mass form leading to (\ref{eq:aux_ineq_mass_stab}).
\end{proof}

\end{document}
